\chapter{Ethernet, Frames and Switching}\label{ch:ethernet}
\bp{1.1}

\begin{outcomes}
By the end of this chapter you will be able to:
\begin{tight}
  \item \textbf{Describe} the Ethernet frame and the purpose of each field you
        will actually meet on a device.
  \item \textbf{Explain} how a switch learns, forwards, floods and ages MAC
        addresses --- the algorithm behind every switching chapter that follows.
  \item \textbf{Read} the MAC address table and use it to prove where a device is
        physically connected.
  \item \textbf{Predict} what a switch does with a unicast, broadcast or unknown
        destination.
\end{tight}
\end{outcomes}

\begin{prereq}
\chref{models} for frames and broadcast domains; \chref{physical} for the link
underneath.
\end{prereq}

\begin{keyterms}
Ethernet $\bullet$ MAC address $\bullet$ OUI $\bullet$ unicast $\bullet$
broadcast $\bullet$ multicast $\bullet$ CAM table $\bullet$ MAC address table
$\bullet$ learning $\bullet$ flooding $\bullet$ filtering $\bullet$ ageing
$\bullet$ unknown unicast $\bullet$ MTU $\bullet$ FCS
\end{keyterms}

\section{The frame}

\begin{definitionbox}[The Ethernet frame, field by field]
\begin{tight}
  \item \textbf{Preamble} --- clock synchronisation. You will never see it.
  \item \textbf{Destination MAC} (6 bytes) --- who this frame is for. The switch
        reads this to decide where to send it.
  \item \textbf{Source MAC} (6 bytes) --- who sent it. The switch reads this
        to \emph{learn}.
  \item \textbf{Type/Length} (2 bytes) --- what is inside. \cmd{0x0800} is IPv4,
        \cmd{0x86DD} is IPv6, \cmd{0x8100} means a VLAN tag follows
        (\chref{trunking}).
  \item \textbf{Payload} --- 46 to 1500 bytes. Below 46 the frame is padded;
        this is why a frame under 64 bytes total is a \term{runt}.
  \item \textbf{FCS} (4 bytes) --- a checksum. If it fails, the frame is counted
        as a CRC error and discarded (\chref{physical}).
\end{tight}
\end{definitionbox}

A \term{MAC address} is 48 bits, written as twelve hex digits. The first half is
the \term{OUI}, assigned to the manufacturer; the second half is assigned by that
manufacturer. Cisco writes them as \cmd{00d0.58a1.0201}; almost everyone else
writes \cmd{00:d0:58:a1:02:01}. They are the same address, and you should be able
to read both.

\begin{conceptbox}[Three kinds of destination]
\begin{tight}
  \item \textbf{Unicast} --- one specific device. The low-order bit of the first
        byte is 0.
  \item \textbf{Broadcast} --- \cmd{ffff.ffff.ffff}, every device in the VLAN.
        This is what defines a broadcast domain.
  \item \textbf{Multicast} --- a group. The low-order bit of the first byte is 1.
        OSPF (\chref{ospfv2}), HSRP and VRRP (\chref{fhrp}) all use multicast.
\end{tight}
\end{conceptbox}

\section{How a switch actually works}

A switch runs one small algorithm, and everything in Part~II is a modification
of it.

\begin{conceptbox}[Learn, then forward or flood]
For every frame that arrives on a port:
\begin{tightnum}
  \item \textbf{Learn.} Record the \emph{source} MAC address against the port it
        arrived on, in this VLAN. Reset the ageing timer if it is already there.
  \item \textbf{Decide on the destination.}
        \begin{tight}
          \item Known unicast $\rightarrow$ \textbf{forward} out that one port.
          \item Unknown unicast $\rightarrow$ \textbf{flood} out every other port
                in the VLAN.
          \item Broadcast or multicast $\rightarrow$ \textbf{flood}.
          \item Destination is on the same port it arrived on $\rightarrow$
                \textbf{filter}, that is, discard it.
        \end{tight}
\end{tightnum}
A switch never modifies the addresses in a frame it forwards, and it never
consults an IP address. Both of those are the router's job.
\end{conceptbox}

\begin{alertbox}[The consequence that matters: unknown unicast flooding]
A switch floods a unicast frame when it does not know the destination. That is
not a fault --- it is how the network keeps working while the table is learned.
But it means a full or wiped MAC table causes traffic to be sprayed to every
port, which is both a performance problem and a security one. It is also the
mechanism a MAC flooding attack abuses, which is why \chref{l2security} exists
and why topic \tp{4.7} names port security.
\end{alertbox}

Entries age out after five minutes of silence by default. A device that has been
quiet longer than that has to be re-learned, and the first frame to it is
flooded.

\begin{verify}{SW1}{show mac address-table}
          Mac Address Table
-------------------------------------------

Vlan    Mac Address       Type        Ports
----    -----------       --------    -----
   1    00d0.58a1.0201    DYNAMIC     Gi0/1
   1    0060.5c2b.19a4    DYNAMIC     Gi0/2
  10    000a.f36e.7712    DYNAMIC     Gi0/5
  10    0002.1751.4c33    STATIC      Gi0/6
  20    00e0.b05f.2210    DYNAMIC     Gi0/8
Total Mac Addresses for this criterion: 5
\end{verify}

This table is the most useful single output on a switch. It answers ``where is
this device physically plugged in?'', and it does so by evidence rather than by
documentation --- which is exactly what topic \tp{2.3} asks you to do when it
says \emph{validate the accuracy of network documentation}.

\begin{config}{a static MAC entry, and a shorter ageing time}
! Pin a known device to a known port. Rarely needed, but it survives
! ageing and cannot be moved by a spoofed frame.
mac address-table static 0002.1751.4c33 vlan 10 interface GigabitEthernet0/6
!
! Age entries out faster than the 300-second default.
mac address-table aging-time 120
\end{config}

\begin{worked}[Following a frame across two switches]
PC-A (\cmd{00d0.58a1.0201}) on \cmd{SW1 Gi0/1} sends to PC-B
(\cmd{00e0.b05f.2210}) on \cmd{SW2 Gi0/3}. Both switches have empty tables.

\begin{tightnum}
  \item The frame arrives at \cmd{SW1 Gi0/1}. \cmd{SW1} \emph{learns}
        \cmd{00d0.58a1.0201} on \cmd{Gi0/1}.
  \item \cmd{SW1} does not know PC-B, so it \emph{floods} --- out its uplink and
        every other access port.
  \item \cmd{SW2} receives it on its uplink and learns \cmd{00d0.58a1.0201} on
        \emph{that uplink port}, not on \cmd{Gi0/1}. A switch only ever knows the
        direction, not the distance.
  \item \cmd{SW2} does not know PC-B either, so it floods. PC-B receives it.
  \item PC-B replies. Now \cmd{SW2} learns PC-B on \cmd{Gi0/3}, and because it
        already knows PC-A, it \emph{forwards} rather than flooding.
  \item \cmd{SW1} learns PC-B on its uplink. From here on the conversation is
        unicast in both directions.
\end{tightnum}
Two floods to start a conversation, then none. That is the whole design.
\end{worked}

\begin{pitfall}
\begin{tight}
  \item \textbf{Expecting a switch to know a distance.} The table says which
        port, never how many hops. Five hundred devices can appear behind one
        uplink port and that is entirely normal.
  \item \textbf{Reading a stale table.} An entry can be up to five minutes old.
        If you have just moved a cable, the table may still show the old port.
  \item \textbf{Confusing the MAC table with the ARP table.} The MAC table maps
        MAC~$\rightarrow$~port and lives on a switch. The ARP table maps
        IP~$\rightarrow$~MAC and lives on any device with an IP address. The exam
        will test that you know which to look at.
  \item \textbf{Forgetting the VLAN column.} The same MAC address can legitimately
        appear in two VLANs. Learning is always per VLAN.
\end{tight}
\end{pitfall}

\begin{breakfix}[Two ports, one MAC address, and a flapping log.]
Users on VLAN 10 report intermittent drops. The log shows:

\begin{verify}{SW1}{show logging | include MACFLAP}
%SW_MATM-4-MACFLAP_NOTIF: Host 000a.f36e.7712 in vlan 10 is flapping
  between port Gi0/5 and port Gi0/12
\end{verify}

\textbf{Diagnosis.} The switch is learning the same source MAC on two ports in
quick succession. A device cannot be in two places, so one of three things is
true: there is a \textbf{physical loop} joining those ports (the same frame is
arriving by two paths --- see \chref{stp}), a device is \textbf{spoofing} that
address, or two devices have been given the \textbf{same MAC} by a bad clone or
a virtualisation misconfiguration.

\textbf{Fix.} Check spanning tree first, because a loop is both the most likely
cause and the most damaging: \cmd{show spanning-tree vlan 10} and look for a port
that should be blocking and is not. If the topology is loop-free, shut
\cmd{Gi0/12}, confirm the flapping stops, and trace what is attached to it.
\end{breakfix}

\begin{lab}{Watching a switch learn}{Packet Tracer}
\textbf{Topology.} Two switches, uplinked; two PCs on \cmd{SW1}, one on
\cmd{SW2}. All in VLAN 1 for now.

\begin{tasks}
  \item Power everything on and immediately run \cmd{show mac address-table} on
        both switches. Record what is there before any traffic.
  \item Ping from PC-A to PC-C. Re-run the command on both switches and record
        which addresses each has learned, and against which port.
  \item Explain, from your own capture, why \cmd{SW2} shows PC-A against its
        uplink rather than against a host port.
  \item Run \cmd{clear mac address-table dynamic}, then ping again and watch the
        table rebuild.
  \item Set \cmd{mac address-table aging-time 15}. Leave a PC idle for thirty
        seconds and confirm its entry disappears.
  \item \textbf{Extension.} Add a second cable between the two switches to create
        a loop. Observe what happens to the MAC table and the log before
        spanning tree settles --- then read \chref{stp} to find out why the
        network survived.
\end{tasks}
\end{lab}

\begin{checkpoint}
\begin{tightnum}
  \item A switch receives a frame whose destination MAC is already in its table,
        against the same port the frame arrived on. What does it do, and what is
        that behaviour called?
  \item Why does the first ping between two devices often show a longer round
        trip than the ones after it?
  \item You need to prove which wall port a particular laptop is using. Which
        table do you consult, on which device, and what do you need to know first?
\end{tightnum}
\end{checkpoint}

\begin{chaptersummary}
\begin{tight}
  \item A switch learns from the \emph{source} address and decides on the
        \emph{destination} address. Everything else follows.
  \item Known unicast is forwarded, unknown unicast and broadcast are flooded,
        and a frame is never sent back out of the port it came in on.
  \item The MAC address table maps MAC to port, per VLAN, and ages entries out
        after five minutes by default.
  \item MAC table is not ARP table. One lives on a switch and maps to ports; the
        other maps IP to MAC and lives on anything with an IP address.
  \item A MAC address flapping between two ports means a loop, a spoof, or a
        duplicate --- and you check for the loop first.
\end{tight}
\end{chaptersummary}

\begin{reviewq}
  \item \textbf{(MCQ)} A switch receives a frame destined for
        \cmd{ffff.ffff.ffff}. What does it do?
        (a)~forwards it to the router (b)~floods it out all ports in the VLAN
        except the ingress port (c)~drops it (d)~sends it to the CPU
  \item \textbf{(MCQ)} Which field does a switch use to populate its MAC address
        table? (a)~destination MAC (b)~source MAC (c)~source IP (d)~Type/Length
  \item \textbf{(Short)} Explain why a switch shows hundreds of MAC addresses
        against a single uplink port, and why this is not a fault.
  \item \textbf{(Short)} Give two distinct causes of a MAC flap message and say
        which you would investigate first, with a reason.
  \item \textbf{(Applied)} A security auditor asks you to prove that no
        unauthorised device is connected to the finance VLAN. Describe, with
        commands, how you would produce that evidence and what its limitations
        are.
\end{reviewq}
